\section{Agreement Evaluation}
\label{sec:agreement_eval}

We study the agreement between different LLM judges and humans on MT-bench and Chatbot Arena datasets. On MT-bench, we also study the agreement among humans. MT-bench represents a small-scale study with controlled human evaluation, while Chatbot Arena represents a larger-scale study with crowdsourced human evaluation in the wild.
\subsection{Setup}

\textbf{MT-bench.}
We generate answers for all 80 questions with 6 models: GPT-4, GPT-3.5, Claude-V1, Vicuna-13B, Alpaca-13B~\cite{alpaca}, and LLaMA-13B~\cite{touvron2023llama}. We then use 2 kinds of judges: LLM judges and 58 expert-level human labelers. The labelers are mostly graduate students so they are considered experts and more skilled than average crowd workers.
We let LLM judges evaluate all pairs and let each human evaluate at least 20 random multi-turn questions. This resulted in around 3K votes for all questions.
The detailed data collection process is in \Cref{sec:data-collection}.

\textbf{Chatbot Arena.}
We randomly sample 3K single-turn votes from 30K arena data, which covers models including GPT-4,  GPT-3.5, Claude, Vicuna-7B/13B, Koala-13B~\cite{koala}, Alpaca-13B, LLaMA-13B, and Dolly-12B. We use two kinds of judges: LLM judges and collected crowd judges (2114 unique IPs).


\textbf{Metrics.}
We define the \textit{agreement} between two types of judges as the probability of randomly selected individuals (but not identical) of each type agreeing on a randomly selected question.
See more explanation in \Cref{subsec:agreement-evaluation}.
\textit{Average win rate} is the average of win rates against all other players.
These metrics can be computed with or without including tie votes.

\subsection{High agreement between GPT-4 and humans}
\label{sec:agreement}

We compute agreement on MT-bench data.
In \Cref{table:agreement_mt_bench}, GPT-4 with both pairwise comparison and single answer grading show very high agreements with human experts.
The agreement under setup S2 (w/o tie) between GPT-4 and humans reaches 85\%, which is even higher than the agreement among humans (81\%). This means GPT-4's judgments closely align with the majority of humans.
We also show that GPT-4's judgments may help humans make better judgments. During our data collection, when a human's choice deviated from GPT-4, we presented GPT-4's judgments to humans and ask if they are reasonable (details in \Cref{subsec:mt-data-collection}). Despite different views, humans deemed GPT-4's judgments reasonable in 75\% of cases and are even willing to change their choices in 34\% of cases.


The data from Arena shows a similar trend, as illustrated by \Cref{table:agreement_arena}. Comparing GPT-4 and other LLM judges, we find they reach a similar non-tie agreement ratio between humans but the number of non-tied votes from GPT-4 is much larger. This means that GPT-4 is more affirmative and less suffered from position bias but other models also perform well when they give an affirmative answer.

In both tables, GPT-4 with single-answer grading matches both pairwise GPT-4 and human preferences very well.
This means GPT-4 has a relatively stable internal rubric.
Although it may sometimes perform slightly worse than pairwise comparison and give more tie votes, it is a more scalable method.


We then perform a breakdown analysis by computing agreement on different model pairs and categories.
We only include non-tied votes.
In \Cref{fig:categorized_agreement}, we observe the agreement between GPT-4 and human progressively increases in line with the performance disparity of the model pairs (i.e., larger win rate difference), from 70\% to nearly 100\%. This suggests that GPT-4 aligns with humans better when significant performance differences exist between the models.

\begin{table}[h]
\centering
\vspace{-1em}
\caption{Agreement between two types of judges on MT-bench. ``G4-Pair'' and ``G4-Single'' denote GPT-4 with pairwise comparison and single-answer grading respectively.
The single-answer grading can be converted into pairwise comparison results for calculating the agreement.
We report two setups: ``S1'' includes non-tie, tie, and inconsistent (due to position bias) votes and counts inconsistent as tie; ``S2'' only includes non-tie votes.
The agreement between two random judges under each setup is denoted as ``R=''.
The top value in each cell is the agreement, and the bottom gray value is \#votes.
}
\footnotesize
\begin{subtable}[t]{0.48\textwidth}
\resizebox{1\columnwidth}{!}{
\begin{tabular}[t]{lrrrr}
\toprule
Setup & \multicolumn{2}{c}{S1 (R = 33\%)} & \multicolumn{2}{c}{S2 (R = 50\%)}  \\
\cmidrule(lr){2-3}\cmidrule(lr){4-5}

Judge & G4-Single & Human & G4-Single & Human \\
\midrule

G4-Pair  & \shortstack{70\% \\ \textcolor{gray}{ 1138 }} & \shortstack{66\% \\ \textcolor{gray}{ 1343 }} & \shortstack{97\% \\ \textcolor{gray}{ 662 }} & \shortstack{\textbf{85\%} \\ \textcolor{gray}{ 859 }}\\
 \midrule
G4-Single & - & \shortstack{60\% \\ \textcolor{gray}{ 1280 }} & - & \shortstack{85\% \\ \textcolor{gray}{ 739 }}\\
 \midrule
Human & - & \shortstack{63\% \\ \textcolor{gray}{ 721 }} & - & \shortstack{\textbf{81\%} \\ \textcolor{gray}{ 479 }}\\
 \bottomrule
\end{tabular}
}
\caption{First Turn}
\label{table:agreement_mt_bench_first_turn}
\end{subtable}
\hfill
\begin{subtable}[t]{0.48\textwidth}
\resizebox{1\columnwidth}{!}{
\centering
\begin{tabular}[t]{lrrrr}
\toprule
Setup & \multicolumn{2}{c}{S1 (R = 33\%)} & \multicolumn{2}{c}{S2 (R = 50\%)} \\
\cmidrule(lr){2-3}\cmidrule(lr){4-5}

Judge & G4-Single & Human & G4-Single & Human \\
\midrule

G4-Pair & \shortstack{70\% \\ \textcolor{gray}{ 1161 }} & \shortstack{66\% \\ \textcolor{gray}{ 1325 }} &  \shortstack{95\% \\ \textcolor{gray}{ 727 }} & \shortstack{\textbf{85\%} \\ \textcolor{gray}{ 864 }} \\
 \midrule
G4-Single & - & \shortstack{59\% \\ \textcolor{gray}{ 1285 }} & - & \shortstack{84\% \\ \textcolor{gray}{ 776 }} \\
 \midrule
Human & - & \shortstack{67\% \\ \textcolor{gray}{ 707 }} & - & \shortstack{\textbf{82\%} \\ \textcolor{gray}{ 474 }} \\
 \midrule
\end{tabular}
}
\caption{Second Turn}
\label{table:table:agreement_mt_bench_first_turn}
\end{subtable}
\vspace{-0em}
\label{table:agreement_mt_bench}
\end{table}










\vspace{2em}
\begin{table}[h]
\centering
\vspace{-1.5em}
\begin{minipage}{0.59\textwidth}
\caption{Agreement between two types of judges on Chatbot Arena.
``G4-S'' denotes GPT-4 with single-answer grading.
``G4'', ``G3.5'' and ``C'' denote GPT-4, GPT-3.5, and Claude with pairwise comparison, respectively.
``H'' denotes human.
The remaining of table follows the same format as \Cref{table:agreement_mt_bench}. 
}
\footnotesize
\resizebox{1\columnwidth}{!}{
\begin{tabular}{lrrrrrrrr}
\toprule
Setup & \multicolumn{4}{c}{S1 (Random = 33\%)} & \multicolumn{4}{c}{S2 (Random = 50\%)} \\
\cmidrule(lr){2-5}\cmidrule(lr){6-9}

Judge & G4-S & G3.5 & C & H & G4-S & G3.5 & C & H \\
\midrule

G4 & \shortstack{72\% \\ \textcolor{gray}{2968}} & \shortstack{66\% \\ \textcolor{gray}{3061}} & \shortstack{66\%
\\ \textcolor{gray}{3062}} & \shortstack{64\% \\ \textcolor{gray}{3066}} & \shortstack{95\% \\ \textcolor{gray}{1967}} & \shortstack{94\% \\ \textcolor{gray}{1788}} & \shortstack{95\% \\ \textcolor{gray}{1712}} & \shortstack{\textbf{87\%} \\ \textcolor{gray}{1944}}\\
 \midrule
G4-S & - & \shortstack{60\% \\ \textcolor{gray}{2964}} & \shortstack{62\% \\ \textcolor{gray}{2964}} & \shortstack{60\% \\ \textcolor{gray}{2968}} & - & \shortstack{89\% \\ \textcolor{gray}{1593}} & \shortstack{91\% \\ \textcolor{gray}{1538}} & \shortstack{85\% \\ \textcolor{gray}{1761}}\\
 \midrule
G3.5 & - & - & \shortstack{68\% \\ \textcolor{gray}{3057}} & \shortstack{54\% \\ \textcolor{gray}{3061}} & - & - & \shortstack{96\% \\ \textcolor{gray}{1497}} & \shortstack{83\% \\ \textcolor{gray}{1567}}\\
 \midrule
C & - & - & - & \shortstack{53\% \\ \textcolor{gray}{3062}} & - & - & - & \shortstack{84\% \\ \textcolor{gray}{1475}}\\
 \bottomrule
\end{tabular}
}
\label{table:agreement_arena}
\end{minipage}
\hfill
\begin{minipage}{0.39\textwidth}
\centering
\includegraphics[width=0.66\textwidth]{figures/categorized_agreement.pdf}
\captionsetup{type=figure}
\caption{Agreement and win rate difference. Each point corresponds to a model pair and counts only the non-tie votes between the two models. The x-axis value is the win rate difference between the two models. The y-axis value is the GPT-4 and human agreement.}
\label{fig:categorized_agreement}
\end{minipage}
\end{table}




\begin{figure}[H]
\centering
\vspace{-2em}
\includegraphics[width=1\textwidth]{figures/mt_win_rate.pdf}
\vspace{-1.5em}
\caption{Average win rate of six models under different judges on MT-bench.}
\vspace{-1em}
\label{fig:mt_bench_win_rate}
\end{figure}

\begin{figure}[H]
\centering
\vspace{-0em}
\includegraphics[width=1\textwidth]{figures/arena_win_rate.pdf}
\vspace{-2em}
\caption{Average win rate of nine models under different judges on Chatbot Arena.}
\vspace{-1em}
\label{fig:arena_win_rate}
\end{figure}


\begin{table}[H]
\centering
\caption{Category-wise win rate of models.}
\label{table:win_rate_category}
\footnotesize
\resizebox{0.95\columnwidth}{!}{
\begin{tabular}{lllllllll}
\toprule
Model & Writing & Roleplay & Reasoning & Math & Coding & Extraction & STEM & Humanities \\
\midrule
GPT-4           & 61.2\%  & 67.9\%  & 49.3\%  & 66.1\%  & 56.3\%  & 66.2\%  & 76.6\%  & 72.2\%  \\
GPT-3.5         & 50.9\%  & 60.6\%  & 32.6\%  & 63.8\%  & 55.0\%  & 48.8\%  & 52.8\%  & 53.8\%  \\
Vicuna-13B      & 39.7\%  & 39.2\%  & 20.1\%  & 18.0\%  & 36.9\%  & 29.2\%  & 47.0\%  & 47.5\%  \\
LLaMA-13B       & 15.1\%  & 15.1\%  & 7.8\%   & 7.5\%   & 2.1\%   & 9.3\%   & 6.8\%   & 10.1\%  \\
\bottomrule
\end{tabular}
}
\end{table}


\subsection{Win rates under different judges}
We plot the average win rate of models under different judges on MT-bench and Chatbot Arena in \Cref{fig:mt_bench_win_rate} and \Cref{fig:arena_win_rate}, respectively. The win rate curves from LLM judges closely match the curves from humans. On MT-bench second turn, proprietary models like Claude and GPT-3.5 are more preferred by the humans compared to the first turn, meaning that a multi-turn benchmark can better differentiate some advanced abilities of models. We also list the per-category win rate of representative models in \Cref{table:win_rate_category} to show how MT-bench differentiates models, in which we see GPT-4 is significantly better than others.
Vicuna-13B is noticeably worse than GPT-3.5/4 in reasoning, math, and coding categories.
Note that in math/coding category, GPT-3.5 and GPT-4 have similar overall win-rate because they both failed to answer some hard questions, but GPT-4 is still significantly better than GPT-3 in the direct pairwise comparison or single-answer grading.
Please see a performance breakdown of MT-bench score for each category in \Cref{subsec:category-wise-scores}.